\subsection{Analisis Kebutuhan Pipeline PGHD}

Fokus: menjawab RM1 di level analisis. Topik yang bisa dibahas:

\begin{enumerate}
    \item Dari literatur PGHD–EHR: PGHD punya karakteristik:
    \begin{itemize}
        \item heterogen (berbagai sumber: smartphone, wearable, self-report),
        
        \item time-series, kontinyu, sering high-frequency.
    \end{itemize}

    \item Dari literatur edge / IoMT:
    \begin{itemize}
        \item tidak efisien mengirim semua raw data ke server/ledger.
        
        \item banyak sistem desain data aggregation pipeline di edge (phone) lalu hanya mengirim fitur/summary/deret waktu terstruktur.
    \end{itemize}

    \item Kebutuhan pipeline di kasusmu:
    \begin{itemize}
        \item mengumpulkan data dari sensor Android + data wearable (simulasi),
        
        \item melakukan transformasi raw → fitur (steps, HR, dsb.),
        
        \item enkripsi di sisi klien,
        
        \item upload ke IPFS, kirim CID + metadata ke IOTA.
    \end{itemize}
\end{enumerate}

Di akhir, kamu bisa merumuskan kebutuhan fungsional pipeline, misalnya:
\begin{enumerate}
    \item P1: pipeline harus mampu meng-handle ingestion data sensor multi-sumber,
    
    \item P2: pipeline harus mendukung enkripsi end-to-end,
    
    \item P3: pipeline harus memetakan data ke schema yang konsisten (mis. FHIR-like).
\end{enumerate}